
\begingroup
\renewcommand*{\mriteabstracttitlefont}{\heiti}
\renewcommand*{\mritekeywordslabelfont}{\song\bfseries}
\begin{abstract}

扩大模型、增加训练数据和改善语料质量都需要消耗算力，其收益又受投入起点与处理费用影响。本文从语料评价出发，建立资源投入与训练损失的关系，求解有限预算下的配置，并利用历史评测分析能力前沿。

\textbf{针对问题一：} 对 272,505 条记录的 22 个质量指标按字段含义统一方向，以分位截断、中位数填补和熵权法合成评分。成对得分差与 Kendall 系数分别衡量样本内分歧和总体排序一致性：冲突比例为 21.9\%，低于独立参照的 28.9\%，但 Kendall 系数仅为 0.076；域内中位数标签只表示相对位置。领域均分以 arxiv 最高、github 最低。对于配比与损失，线性岭保留系数解释，平方根嵌入核岭用于预测及训练配方凸包内推荐。在 1M 检验集，平均逐域 $R^2$ 从 0.6186 提高至 0.9454，综合 RMSE 从 0.5232 降至 0.2256。核岭在固定域和配比上限约束下给出预测加权损失 4.6444，参考配方为 4.8173；这一模型内收益尚需训练实验验证。跨尺度绝对误差较大，故后续仅将配比损失差作为条件修正。质量指数与领域占比线性相关，其重复加入不能识别质量的独立作用。

\textbf{针对问题二：} 采用稳健非线性最小二乘拟合经典标度律，并在四种质量项中选择加性缺口形式。经典模型与广义联合模型的 $R^2$ 分别为 0.9999998、0.9896，跨族和文献记录的 MAPE 则为 8.8\%、8.0\%，显示同源拟合精度不能直接推广到其他来源。模型导数用于计算弹性与局部替代率，完整等损失方程用于较大幅度的投入调整。质量均为 0.6 时，1B 参数、100B tokens 与 12B 参数、300B tokens 两个起点的质量--参数替代率分别为 2.844 和 76.838 B/单位质量。领域系数的中心化相似度描述作用方向，核岭相对参考配方的损失差另作为配比修正；质量效应与跨尺度迁移仍受半合成实验和尺度假设限制。

\textbf{针对问题三：} 将基础训练、质量处理及注意力费用计入预算，以广义模型的预测损失为目标，采用 SLSQP 求解。幂函数型质量成本、参考配比及 2048 tokens 上下文下，$10^{19}$ FLOPs 对应 0.225B 参数、6.9B tokens，质量维持在基线 0.549；$10^{22}$、$10^{24}$ FLOPs 下质量均为 1，参数与数据配置分别为 6.091B、229.4B tokens 和 44.939B、3417.8B tokens。扫描显示维持基线、增加质量和达到上界三个阶段，首次投入与饱和的有效网格点为 $1.58\times10^{19}$、$7.08\times10^{19}$ FLOPs。成本式给出 30,000 tokens 的注意力与基础训练等费用长度。固定 $10^{22}$ FLOPs，上下文增至 131072 tokens 时参数量降至 2.705B，损失约增加 6.0\%。不同质量成本与尺度映射改变投入门槛，配比收益的预算等价也依赖迁移幅度。

\textbf{针对问题四：} 从许可与字段筛选后的 2,346 条记录出发，以六项评测均值描述能力。损失与得分桥接的 $R^2$ 仅为 0.26，因此只解释关联方向。合并横截面回归中，参数量每扩大十倍对应 13.84 分差异，控制规模后的年度组差为 4.84 分，后者不表示同一模型的年增长。将前沿评分与相应规模配对后，2019—2025 年窗口的规模关联项占 51.4\%，剩余项占 48.6\%；2024 年 6—9 月窗口的剩余项占 95.4\%，其中仍含样本构成等未观测影响。以 2025 年 3 月为预测基点，logistic 的 12、24 个月条件预测为 56.3、56.6 分，曲线重拟合的 5\%--95\% 分位范围为 53.6--64.5、53.8--66.5。时间留出未显示稳定预测优势，窗口删减也会改变结果。增长参数减半时两点为 55.6、56.3 分；另一结构情景将设定年增量减半时为 53.4、55.8 分，两类情景的基期水平不同。C8 子任务与 C4 资源记录进一步用于比较能力短板和数据覆盖。

各项配置和外推均对应明确的成本、评分与观测窗口。质量处理应随预算和边际费用安排，模型扩容则须结合可用数据判断；跨来源标定不足的部分保留为条件分析。

\keywords{大语言模型；标度律；数据质量；资源配置优化；能力前沿预测}
\end{abstract}
\endgroup
